\documentclass{article}
% === 论文样式包 ===
\usepackage{PRIMEarxiv}
\usepackage{natbib}
\usepackage{CJKutf8}
\usepackage{CJK}
\usepackage[utf8]{inputenc} % 允许utf-8输入
\usepackage[T1]{fontenc}    % 使用8位T1字体
\usepackage{hyperref}       % 超链接
\usepackage{url}            % URL排版
\usepackage{booktabs}       % 专业表格
\usepackage{amsfonts}       % 数学黑板符号
\usepackage{nicefrac}       % 紧凑分数符号
\usepackage{microtype}      % 微排版
\usepackage{lipsum}
\usepackage{graphicx}
\usepackage{amsmath}
\usepackage{pifont}
\usepackage[inline]{enumitem}
\usepackage{enumitem}
\usepackage{boldline}
\usepackage{hhline}
\usepackage{multirow}
\usepackage{xcolor}
\usepackage{fontawesome5}
\usepackage{tablefootnote}
\usepackage{listings}
\usepackage{booktabs,threeparttable,siunitx,xcolor,pgfplots}

% === 自定义盒子样式 ===
\usepackage{xspace}
\usepackage[most]{tcolorbox}

% 系统盒子
\newtcolorbox{sysbox}[1][]{
  colback=gray!6,
  colframe=black!20,
  boxrule=0.5pt,
  arc=1.5mm,
  left=1mm,right=1mm,top=1mm,bottom=1mm,
  #1
}

% 提示词盒子
\newtcolorbox{promptbox}[2][]{
  enhanced,
  breakable,
  colback=black!2,
  colframe=black!70,
  boxrule=0.7pt,
  arc=2pt,
  left=6pt,right=6pt,top=6pt,bottom=6pt,
  title={#2},
  fonttitle=\bfseries,
  #1
}

% 概览盒子
\newtcolorbox{overviewbox}{
  colback=gray!6,
  colframe=black!20,
  boxrule=0.5pt,
  arc=1.2mm,
  left=1.2mm,right=1.2mm,top=1mm,bottom=1mm
}

% 文件布局盒子
\newtcolorbox{layoutbox}[1]{
  colback=black!2,
  colframe=black!18,
  boxrule=0.45pt,
  arc=1mm,
  left=1.5mm,right=1.5mm,top=1mm,bottom=1mm,
  title=\textbf{#1},
  fonttitle=\small,
  coltitle=black
}

% === 颜色定义 ===
\definecolor{skillblue}{HTML}{2F5D8A}
\definecolor{skilllight}{HTML}{F4F8FC}
\definecolor{skillborder}{HTML}{A9C3DE}
\definecolor{tagbg}{HTML}{EAF2FB}
\definecolor{dagpurple}{HTML}{6B5B95}
\definecolor{agentgreen}{HTML}{4CAF50}
\definecolor{planorange}{HTML}{FF9800}

% 标签样式
\newtcbox{\skilltag}{
  on line,
  boxrule=0pt,
  arc=2mm,
  colback=tagbg,
  colframe=tagbg,
  left=2mm,right=2mm,top=0.6mm,bottom=0.6mm
}

% 列表样式优化
\setlist[itemize]{leftmargin=1.5em, topsep=2pt, itemsep=2pt, parsep=0pt}
\setlist[enumerate]{leftmargin=1.8em, topsep=2pt, itemsep=2pt, parsep=0pt}
\setlist[description]{leftmargin=1.6em, style=nextline, topsep=2pt, itemsep=3pt}

% 图标命令
\newcommand{\hficon}{\raisebox{-0.2ex}{\includegraphics[height=1.8ex]{figure/hug.png}}}
\newcommand{\lockicon}{\raisebox{-0.2ex}{\includegraphics[height=1.8ex]{figure/lock.png}}}
\newcommand{\brainicon}{\raisebox{-0.2ex}{\includegraphics[height=1.8ex]{figure/brain.png}}}

\graphicspath{{images/}}
\usepackage{amsthm}
\theoremstyle{definition}
\newtheorem{definition}{Definition}

\usepackage{verbatim}
\usepackage[most]{tcolorbox}
\tcbuselibrary{breakable, listings}

% Finding Box定义
\newtcolorbox{findingbox}{
  colback=black!5!white,
  colframe=black!75!black,
  arc=3mm,
  boxrule=0.5pt,
  left=4mm, right=4mm, top=3mm, bottom=3mm
}

% textBox定义
\newtcolorbox{textBox}[1][]{
  colback=gray!5,
  colframe=gray!60!black,
  fonttitle=\bfseries,
  colbacktitle=gray!85!black,
  title=#1,
  breakable
}

% === 论文标题 ===
\title{ElianDataAgent: A Plan-Execute-Replan Multi-Agent Framework with DAG-Enhanced Parallel Execution for Data Analysis Tasks}

% === 作者信息 ===
\author{
  \textbf{Zimo Li}$^{1}$\footnotemark[1] \\
  $^1$ School of Mathematics and Statistics, Nanjing University of Information Science and Technology \\
  \texttt{\{lizimo\}@nuist.edu.cn}, \\
  \textcolor{red}{\url{https://github.com/2Elian/elian_data_agent}}
}

\begin{document}

\renewcommand{\thefootnote}{\fnsymbol{footnote}}
\maketitle

% % === 摘要部分 ===

% 现实世界中的数据分析任务通常需要复杂的多步骤推理、工具编排以及对中间结果的动态调整。

% 现有的单智能体方法在任务分解、并行执行和故障恢复方面存在困难，导致工作流效率低下且可扩展性有限。

% 我们提出了 ElianDataAgent，

% 这是一个多智能体框架，它将计划-执行-重计划范式与有向无环图 (DAG) 增强的并行执行相结合，以应对复杂的数据分析任务。

% ElianDataAgent 采用三个专用智能体：一个规划器，用于将任务分解为具有显式依赖关系的结构化执行计划；一个执行器，使用 ReAct 模式执行分步推理；以及一个重规划器，用于根据执行结果动态调整计划。

% 该系统通过 DAG 调度器进一步提高执行效率，该调度器支持三种执行策略：顺序执行、并行执行和混合执行，从而实现依赖关系感知的并发任务处理。

% 实验评估表明，ElianDataAgent 能够有效处理各种数据分析场景，包括 CSV 探索、SQL 查询和多源数据集成，并凭借自适应重规划能力实现稳健的任务完成。该框架为构建智能数据分析代理提供了一个实用且可扩展的基础，这些代理能够将战略规划与高效的并行执行相结合。

\begin{abstract}
% 中文注释：摘要介绍系统核心思想和主要贡献
Data analysis tasks in real-world scenarios often require complex multi-step reasoning,
tool orchestration, and dynamic adaptation to intermediate results.
 Existing single-agent approaches struggle with task decomposition, 
 parallel execution, and failure recovery, leading to inefficient workflows and limited scalability. 
 We present ElianDataAgent, 
 a multi-agent framework that integrates the Plan-Execute-Replan paradigm with Directed Acyclic 
 Graph enhanced parallel execution for sophisticated data analysis tasks.
ElianDataAgent employs three specialized agents: a Planner that decomposes tasks into structured execution 
plans with explicit dependencies, an Executor that performs step-wise reasoning using the ReAct 
pattern, and a Replanner that dynamically adjusts plans based on execution outcomes. 
The system further enhances execution efficiency through a DAG scheduler that supports three execution 
strategies: Sequential, Parallel, and Hybrid that enabling dependency-aware concurrent task processing.
Experimental evaluation demonstrates that ElianDataAgent effectively handles diverse data ana
lysis scenarios including CSV exploration, SQL querying, and multi-source data integration, achiev
ing robust task completion with adaptive replanning capabilities. The framework provides a practical
 and scalable foundation for building intelligent data analysis agents that combine strategic 
 planning with efficient parallel execution.
\end{abstract}

% === 关键词 ===
\keywords{Multi-Agent Systems \and Plan-Execute-Replan \and DAG Execution \and Data Analysis \and ReAct Pattern}

% === 第一章：引言 ===
\section{Introduction}

% 中文注释：引言背景部分，介绍数据分析任务的复杂性和现有系统的局限性
Large language models (LLMs) have demonstrated remarkable capabilities in 
natural language understanding, code generation, and tool use 
\cite{brown2020language,touvron2023llama}. 
These advances have enabled a new generation of 
intelligent agents that can perform complex reasoning tasks, 
interact with external tools, and solve multi-step problems autonomously 
\cite{yao2022react,schick2023toolformer,patil2024gorilla}. 
However, real-world data analysis tasks present unique challenges
 that exceed the capabilities of single-agent systems.

% 中文注释：具体挑战列举
Data analysis workflows typically involve multiple interconnected operations: exploring file structures, reading diverse data formats (CSV, JSON, SQLite), formulating and executing SQL queries, performing statistical computations, and synthesizing final answers. These tasks exhibit several characteristics that complicate single-agent approaches:

\begin{itemize}[leftmargin=*, align=left]
    \item \textbf{Complex task decomposition:} High-level queries like ``analyze the sales trends across multiple tables'' require sophisticated planning to identify relevant data sources, determine appropriate analysis methods, and sequence operations correctly.
    \item \textbf{Dependency management:} Certain operations must precede others (e.g., schema inspection before query formulation), while independent operations could potentially execute concurrently for efficiency.
    \item \textbf{Dynamic adaptation:} Intermediate results may reveal unexpected patterns or data quality issues, requiring plan modifications without complete restarts.
    \item \textbf{Error recovery:} Failed operations (e.g., malformed queries, missing columns) need graceful handling with alternative approaches rather than task failure.
\end{itemize}

% 中文注释：现有方法的不足
Existing approaches to LLM-based data analysis predominantly employ single-agent architectures with ReAct-style reasoning loops \cite{yao2022react}. While effective for straightforward tasks, these systems face fundamental limitations: they cannot parallelize independent operations, struggle with complex dependency management, and lack systematic mechanisms for plan revision. Memory-enhanced agents \cite{zhong2024memorybank,packer2023memgpt} improve context retention but do not address structural planning challenges. Specialized tool-learning methods \cite{schick2023toolformer,tang2023toolalpaca,qintoolllm} enhance individual tool use but lack orchestration frameworks for multi-tool coordination.

% 中文注释：本论文的贡献
In this paper, we present \textbf{ElianDataAgent}, a multi-agent framework that addresses these challenges through two integrated design principles: the Plan-Execute-Replan (PER) paradigm and DAG-enhanced parallel execution. The framework draws inspiration from established agent architectures---specifically the Eino Agent Development Kit's session-based agent design and Shannon's DAG workflow patterns---while introducing novel adaptations for data analysis domains.

% 中文注释：系统核心组件介绍
The ElianDataAgent architecture comprises three specialized agents operating within a shared Session context:

\begin{itemize}[leftmargin=*, align=left]
    \item \textbf{Planner Agent:} Decomposes user queries into structured execution plans with explicit step dependencies, supporting Sequential, DAG, and Hybrid execution modes.
    \item \textbf{Executor Agent:} Executes individual plan steps through ReAct-style reasoning (Think-Act-Observe loops) with access to a validated tool registry for data manipulation.
    \item \textbf{Replanner Agent:} Reviews execution outcomes and dynamically adjusts plans---adding, modifying, or removing steps based on intermediate results and error conditions.
\end{itemize}

% 中文注释：DAG执行引擎特点
Beyond agent specialization, ElianDataAgent introduces a DAG execution engine that enables dependency-aware parallel execution. Using Kahn's algorithm for topological sorting and a three-mode execution strategy (Sequential, Parallel, Hybrid), the system maximizes throughput while respecting operation dependencies. A shared Workspace component manages cold-start schema discovery, file caching, and knowledge integration, providing all agents with consistent context about available data sources.

% 中文注释：主要贡献总结
The main contributions of this paper are:
\begin{itemize}[leftmargin=*, align=left]
    \item We propose a multi-agent architecture that combines Plan-Execute-Replan orchestration with DAG-based parallel execution, enabling both strategic planning and efficient task processing for data analysis scenarios.
    \item We design a comprehensive type system for inter-agent communication, including Plan/PlanStep structures, Session state management, and History tracking mechanisms.
    \item We implement three execution strategies (Sequential, Parallel, Hybrid) with ThreadPoolExecutor-based concurrency control and incremental dependency checking.
    \item We develop a secure tool registry with eight validated data manipulation tools featuring path validation, SQL injection prevention, and execution timeouts.
    \item We provide a modular, extensible implementation compatible with OpenAI 0.x and 1.x APIs, suitable for both research experimentation and practical deployment.
\end{itemize}


% === 第二章：相关工作 ===
\section{Related Work}

% 中文注释：相关工作综述，分为四个研究方向
We organize related work into four research threads: multi-agent orchestration patterns, reasoning-acting frameworks, DAG-based task execution, and tool-augmented language models. ElianDataAgent synthesizes insights from all four directions while introducing novel adaptations for data analysis tasks.

% === 2.1 多智能体编排模式 ===
\subsection{Multi-Agent Orchestration Patterns}

% 中文注释：多智能体系统的发展历程
Multi-agent systems for LLM orchestration have emerged as a powerful paradigm for decomposing complex tasks into specialized subtasks. Early frameworks like AutoGen \cite{wu2023autogen} demonstrated conversational agent collaboration through message passing, while BabyAGI \cite{babyagi2023} introduced task-driven execution loops with vector-based task storage. These systems highlighted the value of agent specialization but lacked structured dependency management.

% 中文注释：Plan-Execute-Replan模式介绍
The Plan-Execute-Replan (PER) pattern, formalized in systems like Plan-and-Execute agents \cite{plansolve2024}, provides a principled approach to task decomposition and adaptive execution. In PER architectures, a planner generates structured execution plans, an executor processes individual steps, and a replanner dynamically adjusts plans based on outcomes. ElianDataAgent adopts this three-agent structure but extends it with DAG-based execution ordering and dependency-aware parallelization.

% 中文注释：会话状态管理相关研究
Session-based state management, inspired by frameworks like Eino ADK, enables consistent information sharing across agents through centralized KV stores. ElianDataAgent implements a Session component that maintains execution state, intermediate results (as topics), and plan history, providing a unified context for all agents.

% === 2.2 推理-行动框架 ===
\subsection{Reasoning-Acting Frameworks}

% 中文注释：ReAct模式的核心思想
The ReAct (Reasoning + Acting) framework \cite{yao2022react} introduced interleaved thinking and acting as a core paradigm for tool-using agents. In ReAct loops, agents alternate between generating thoughts (reasoning about current state), actions (tool invocations), and observations (tool outputs), enabling transparent reasoning traces and iterative refinement. This pattern has been widely adopted in agent benchmarks including WebShop \cite{yao2022webshop}, WebArena \cite{zhou2023webarena}, and ALFWorld \cite{shridharalfworld}.

% 中文注释：ElianDataAgent对ReAct的应用
ElianDataAgent integrates ReAct-style reasoning within its Executor agent, enabling step-wise tool use with observable reasoning traces. However, unlike single-agent ReAct systems, ElianDataAgent confines ReAct loops to individual step execution, delegating plan-level decisions to the Planner and Replanner agents. This separation enables both fine-grained tool reasoning and macro-level plan management.

% 中文注释：思维链相关研究
Chain-of-thought prompting \cite{wei2022chain} and its extensions (Tree of Thoughts \cite{yao2023tree}, Graph of Thoughts \cite{besta2024graph}) improve reasoning quality but operate within single-agent contexts. ElianDataAgent's multi-agent structure complements these approaches by distributing reasoning responsibilities across specialized agents.

% === 2.3 DAG任务执行 ===
\subsection{DAG-Based Task Execution}

% 中文注释：DAG执行模型的理论基础
Directed Acyclic Graphs (DAGs) provide a natural representation for tasks with dependencies, enabling both topological ordering for sequential execution and parallel execution of independent nodes. Work on DAG-based workflow systems---including Apache Airflow, Luigi, and Shannon's DAG patterns \cite{shannon2024dag}---has established principles for dependency management, cycle detection, and concurrent execution.

% 中文注释：拓扑排序算法介绍
Kahn's algorithm \cite{kahn1962topological} remains a standard method for topological sorting with efficient cycle detection. ElianDataAgent employs Kahn's algorithm in its DAGGraph component, computing execution layers that reveal which nodes can execute concurrently. The DAGExecutor implements three execution strategies:
\begin{itemize}[leftmargin=1.5em]
    \item \textbf{Sequential:} Execute all nodes in topological order, suitable for tasks with tight dependencies.
    \item \textbf{Parallel:} Execute all nodes concurrently, applicable when no dependencies exist.
    \item \textbf{Hybrid:} Execute independent nodes concurrently while respecting dependencies, achieving optimal throughput.
\end{itemize}

% 中文注释：并发执行控制机制
Concurrency control through ThreadPoolExecutor with semaphore-based throttling prevents resource exhaustion while enabling genuine parallelism. Incremental dependency checking with configurable timeouts handles long-running dependencies gracefully.

% === 2.4 工具增强语言模型 ===
\subsection{Tool-Augmented Language Models}

% 中文注释：工具学习相关研究
Tool use has become a core capability for LLM agents, enabling interaction with external APIs, databases, and computational resources. Toolformer \cite{schick2023toolformer} demonstrated self-supervised tool learning, while Gorilla \cite{patil2024gorilla} trained models specifically for API invocation. Tool-oriented benchmarks (API-Bank \cite{li2023api}, ToolBench \cite{qintoolllm}, ToolAlpaca \cite{tang2023toolalpaca}) evaluate tool use competence across diverse scenarios.

% 中文注释：ElianDataAgent的工具注册表设计
ElianDataAgent implements a secure tool registry with eight validated data manipulation tools: file listing, CSV reading, JSON parsing, document inspection, SQLite schema queries, SQL execution, Python code execution, and answer submission. Each tool includes input validation (path traversal prevention, SQL injection protection) and timeout constraints, ensuring safe operation in data analysis contexts.


% === 第三章：方法 ===
\section{Method}

% 中文注释：方法章节概述
We propose a multi-agent framework for data analysis that combines Plan-Execute-Replan orchestration with DAG-enhanced parallel execution. The core design separates strategic planning, step-wise execution, and dynamic replanning across three specialized agents, while a DAG scheduler enables efficient concurrent task processing.

% === 3.1 问题定义 ===
\subsection{Problem Definition}

% 中文注释：问题形式化定义
Consider a data analysis task specified by a natural language query $q$ over a collection of data sources $\mathcal{D} = \{d_1, d_2, \dots, d_n\}$ (CSV files, JSON files, SQLite databases). The task requires producing an answer $a$ that addresses $q$ through a sequence of operations including:
\begin{itemize}[leftmargin=1.5em]
    \item Data discovery: identifying relevant data sources within $\mathcal{D}$.
    \item Schema exploration: understanding table structures, column names, and data types.
    \item Query formulation: constructing appropriate SQL queries or analysis code.
    \item Execution: running queries, computing statistics, or transforming data.
    \item Synthesis: aggregating results into a coherent answer.
\end{itemize}

% 中文注释：执行计划的形式化表示
Let $\pi = (s_1, s_2, \dots, s_m)$ denote an execution plan where each step $s_i$ specifies an operation. Steps may have dependencies: $s_i \prec s_j$ indicates $s_i$ must complete before $s_j$ can begin. The dependency structure forms a DAG $\mathcal{G} = (V, E)$ where $V = \{s_1, \dots, s_m\}$ and $E$ encodes dependency relations.

% 中文注释：系统目标陈述
Our goal is to design a system that:
\begin{enumerate}[leftmargin=1.5em]
    \item Generates appropriate execution plans $\pi$ from queries $q$.
    \item Executes plans efficiently, maximizing parallelism while respecting dependencies.
    \item Adapts plans dynamically when intermediate results reveal unexpected conditions.
    \item Produces accurate answers through systematic result synthesis.
\end{enumerate}

% === 3.2 系统架构 ===
\subsection{System Architecture}

% 中文注释：架构总览
ElianDataAgent implements a four-layer architecture as shown in Figure~\ref{fig:framework}:

\begin{overviewbox}
\noindent\textbf{Architecture Overview.}
ElianDataAgent comprises four layers: (1) a Core layer defining types and shared workspace, (2) an Agent layer with Planner, Executor, and Replanner, (3) a DAG layer for dependency management and parallel execution, and (4) a Tool layer providing validated data manipulation capabilities. All agents operate within a shared Session context, enabling consistent state management across orchestration phases.
\end{overviewbox}

% === 3.2.1 核心层 ===
\subsubsection{Core Layer}

% 中文注释：核心层的两个组件
The Core layer provides foundational abstractions used across all system components:

\paragraph{Type System.}
% 中文注释：类型系统的定义
We define a comprehensive type system for inter-agent communication:
\begin{itemize}[leftmargin=1.5em]
    \item \textbf{Plan:} A structured execution plan containing a sequence of PlanSteps with dependency declarations and execution mode preferences.
    \item \textbf{PlanStep:} An individual operation specification including a unique ID, description, dependencies (\texttt{depends\_on}), and output topics (\texttt{produces}).
    \item \textbf{Session:} A shared state container maintaining the current Plan, execution results (as topics), and metadata. Session enables cross-agent information sharing through a KV store pattern.
    \item \textbf{History:} An event tracking mechanism that records agent interactions, tool invocations, and state transitions as AgentEvents.
\end{itemize}

% 中文注释：类型定义的形式化表示
Formally, a Plan is defined as:
\[
\pi = (\text{id}, \text{steps}, \text{plan\_type}, \text{status}, \text{context})
\]
where $\text{steps} = \{s_1, \dots, s_m\}$, each $s_i = (\text{id}, \text{desc}, \text{depends\_on}, \text{produces}, \text{status})$.

A Session maintains:
\[
\mathcal{S} = (\text{plan}, \text{topics}, \text{metadata})
\]
where $\text{topics}$ is a KV store mapping topic names to execution results, enabling data flow between dependent steps.

\paragraph{Workspace.}
% 中文注释：Workspace组件的功能
The Workspace component manages data source discovery and caching:
\begin{itemize}[leftmargin=1.5em]
    \item \textbf{Cold Start:} A single comprehensive scan of the data directory, building a SchemaMap that records file types, table structures, and column information.
    \item \textbf{File Cache:} In-memory caching of read file contents to avoid redundant I/O operations.
    \item \textbf{Knowledge Integration:} Optional loading of domain knowledge (e.g., \texttt{knowledge.md}) for specialized analysis tasks.
\end{itemize}

% === 3.2.2 智能体层 ===
\subsubsection{Agent Layer}

% 中文注释：三个专门化智能体的设计
The Agent layer implements three specialized agents with distinct responsibilities:

\paragraph{Planner Agent.}
% 中文注释：Planner的功能和工作流程
The Planner decomposes natural language queries into structured execution plans. Given a query $q$ and workspace schema information, the Planner produces:
\begin{enumerate}[leftmargin=1.5em]
    \item A sequence of PlanSteps with clear operation descriptions.
    \item Dependency declarations between steps (using \texttt{depends\_on}).
    \item Output topic specifications (using \texttt{produces}).
    \item A recommended execution mode (Sequential, DAG, or Hybrid).
\end{enumerate}

The Planner's decision process considers:
\begin{itemize}[leftmargin=1.5em]
    \item Available data sources and their relationships.
    \item Query complexity and required analysis depth.
    \item Potential parallelization opportunities.
    \item Domain knowledge for specialized reasoning.
\end{itemize}

\paragraph{Executor Agent.}
% 中文注释：Executor的ReAct执行循环
The Executor performs individual PlanSteps using ReAct-style reasoning. For each step $s_i$, the Executor enters a Think-Act-Observe loop:

\begin{enumerate}[leftmargin=1.5em]
    \item \textbf{Think:} Analyze current state, identify required tools, and formulate action plans.
    \item \textbf{Act:} Invoke appropriate tools from the registry with validated inputs.
    \item \textbf{Observe:} Process tool outputs, update internal state, and determine next action.
    \item \textbf{Repeat:} Continue until step completion or maximum iteration threshold.
\end{enumerate}

The Executor maintains per-step iteration limits to prevent runaway loops and tracks all tool invocations for History recording. Upon completion, results are stored in Session as topics for downstream steps.

\paragraph{Replanner Agent.}
% 中文注释：Replanner的决策机制
The Replanner reviews execution outcomes and determines next actions. After each execution phase, the Replanner evaluates:
\begin{itemize}[leftmargin=1.5em]
    \item Step completion status (success, failure, partial).
    \item Result quality and relevance to the original query.
    \item Remaining plan steps and their feasibility.
    \item Potential need for additional operations not in the original plan.
\end{itemize}

The Replanner outputs one of four decisions:
\begin{itemize}[leftmargin=1.5em]
    \item \textbf{Continue:} Proceed with the next execution phase.
    \item \textbf{Replan:} Modify the Plan (add/remove/modify steps) and continue.
    \item \textbf{Complete:} Synthesize final answer from accumulated results.
    \item \textbf{Fail:} Report task failure with diagnostic information.
\end{itemize}

% === 3.2.3 DAG层 ===
\subsubsection{DAG Layer}

% 中文注释：DAG层的三个组件
The DAG layer implements dependency-aware execution through three components:

\paragraph{DAGGraph.}
% 中文注释：DAG数据结构
The DAGGraph data structure manages nodes and edges:
\begin{itemize}[leftmargin=1.5em]
    \item Node representation: Each PlanStep maps to a DAGNode with adjacency information.
    \item Edge management: Dependency relations encoded as directed edges.
    \item Topological sort: Kahn's algorithm computes execution ordering and detects cycles.
    \item Layer computation: Groups nodes by execution depth, revealing parallelizable sets.
\end{itemize}

Kahn's algorithm proceeds as follows:
\begin{enumerate}[leftmargin=1.5em]
    \item Initialize queue with nodes having zero incoming edges.
    \item For each node: decrement indegree of dependents; enqueue when indegree reaches zero.
    \item Track execution order; if unprocessed nodes remain, a cycle exists.
\end{enumerate}

\paragraph{DAGScheduler.}
% 中文注释：调度策略选择
The DAGScheduler selects execution strategies based on dependency analysis:
\begin{itemize}[leftmargin=1.5em]
    \item \textbf{No dependencies:} Select Parallel strategy for maximum concurrency.
    \item \textbf{Tight dependencies:} Select Sequential strategy for ordered execution.
    \item \textbf{Mixed dependencies:} Select Hybrid strategy for dependency-aware parallelism.
\end{itemize}

The scheduler also converts Plans to DAGGraphs, mapping PlanStep dependencies to DAG edges.

\paragraph{DAGExecutor.}
% 中文注释：执行引擎的实现
The DAGExecutor implements the selected strategy using ThreadPoolExecutor:
\begin{itemize}[leftmargin=1.5em]
    \item \textbf{Parallel mode:} Submit all tasks concurrently; wait for completion.
    \item \textbf{Sequential mode:} Execute in topological order; single-threaded processing.
    \item \textbf{Hybrid mode:} Execute by layers; within each layer, parallelize independent tasks.
\end{itemize}

Dependency handling in Hybrid mode uses incremental checking:
\begin{enumerate}[leftmargin=1.5em]
    \item Poll dependent step status at configurable intervals (e.g., 5 seconds).
    \item Proceed when dependencies satisfy completion criteria.
    \item Timeout handling for long-running dependencies.
\end{enumerate}

% === 3.2.4 工具层 ===
\subsubsection{Tool Layer}

% 中文注释：工具注册表的设计
The Tool layer provides validated data manipulation capabilities through a structured registry:

\begin{table}[t]
\centering
\caption{Tool registry components with validation mechanisms.}
\label{tab:tools}
\begin{tabular}{lll}
\toprule
Tool Name & Function & Validation \\
\midrule
list\_context & Enumerate files/directories & Path traversal check \\
read\_csv & Parse CSV with preview & File existence, size limit \\
read\_json & Parse JSON with preview & Format validation \\
read\_doc & Read text documents & Encoding detection \\
inspect\_sqlite\_schema & Query table structures & Database validation \\
execute\_context\_sql & Execute read-only SQL & Injection prevention \\
execute\_python & Run Python code & Timeout (30s), sandbox \\
answer & Submit final result & Terminal action \\
\bottomrule
\end{tabular}
\end{table}

Each tool implements input validation, execution constraints, and error handling. The SQL execution tool restricts operations to SELECT statements, preventing destructive queries. Python execution includes a 30-second timeout to prevent runaway computations.

% === 3.3 工作流程编排 ===
\subsection{Workflow Orchestration}

% 中文注释：主工作流的四个阶段
The PlanExecuteReplanWorkflow orchestrates the complete execution cycle across four phases:

\begin{enumerate}[leftmargin=*, label=\textbf{Phase \arabic*.}]
    \item \textbf{Workspace Initialization.}
    Cold-start scan builds SchemaMap; caches file contents; loads optional knowledge.

    \item \textbf{Planning.}
    Planner generates initial Plan from query and workspace context; DAGScheduler converts Plan to DAGGraph.

    \item \textbf{Execution Loop.}
    \begin{enumerate}[leftmargin=1.5em]
        \item DAGExecutor identifies ready steps (completed dependencies).
        \item Executor processes each ready step through ReAct loops.
        \item Results stored in Session as topics.
        \item Replanner evaluates outcomes and decides next action.
        \item Loop continues until Complete or Fail decision.
    \end{enumerate}

    \item \textbf{Answer Synthesis.}
    Upon Complete decision, synthesize final answer from Session topics and History events.
\end{enumerate}

% 中文注释：执行模式的区分
The workflow supports two execution modes:
\begin{itemize}[leftmargin=1.5em]
    \item \textbf{Sequential Mode:} Traditional Plan-Execute-Replan without DAG enhancement; suitable for simple tasks or when DAG conversion fails.
    \item \textbf{DAG Mode:} Full DAG-enhanced execution with dependency-aware parallelism; optimal for complex multi-source tasks.
\end{itemize}


% === 第四章：系统实现 ===
\section{System Implementation}

% 中文注释：系统实现章节概述
ElianDataAgent is implemented as a modular Python system with clear component boundaries and extensibility hooks. We describe the implementation structure, key design decisions, and integration mechanisms.

% === 4.1 代码结构 ===
\subsection{Code Structure}

% 中文注释：项目文件组织
The project follows a layered directory structure matching the architectural design:

\begin{layoutbox}{ElianDataAgent Directory Layout}
\small
\begin{tabular}{p{0.53\linewidth} p{0.41\linewidth}}
\ttfamily
elian\_data\_agent/ & Root directory for the multi-agent system. \\
\ttfamily\hspace*{1em}core/ & Core layer: types and workspace. \\
\ttfamily\hspace*{2em}types.py & Type definitions: Plan, Session, History. \\
\ttfamily\hspace*{2em}workspace.py & Workspace: cold-start, schema, cache. \\
\ttfamily\hspace*{1em}agents/ & Agent layer: specialized agents. \\
\ttfamily\hspace*{2em}base.py & Agent interface and ChatModelAgent base. \\
\ttfamily\hspace*{2em}planner.py & Planner: task decomposition. \\
\ttfamily\hspace*{2em}executor.py & Executor: ReAct execution. \\
\ttfamily\hspace*{2em}replanner.py & Replanner: plan adaptation. \\
\ttfamily\hspace*{2em}model\_adapter.py & OpenAI API compatibility layer. \\
\ttfamily\hspace*{1em}dag/ & DAG layer: graph and execution. \\
\ttfamily\hspace*{2em}graph.py & DAGGraph: Kahn's algorithm. \\
\ttfamily\hspace*{2em}executor.py & DAGExecutor: parallel strategies. \\
\ttfamily\hspace*{2em}scheduler.py & DAGScheduler: strategy selection. \\
\ttfamily\hspace*{1em}orchestration/ & Workflow orchestration. \\
\ttfamily\hspace*{2em}workflow.py & PlanExecuteReplanWorkflow. \\
\ttfamily\hspace*{1em}tools/ & Tool layer: validated tools. \\
\ttfamily\hspace*{2em}registry.py & Tool registry and validation. \\
\ttfamily\hspace*{1em}runner.py & Entry point and configuration. \\
\end{tabular}
\end{layoutbox}

% === 4.2 模型适配 ===
\subsection{Model Adapter}

% 中文注释：API兼容性设计
ElianDataAgent supports both OpenAI API versions (0.x and 1.x) through a ModelAdapter compatibility layer. The adapter:
\begin{itemize}[leftmargin=1.5em]
    \item Auto-detects installed OpenAI version.
    \item Normalizes message formats across versions.
    \item Handles API base URL and key configuration.
    \item Provides consistent interface for all agents.
\end{itemize}

This design enables deployment across diverse environments without version-specific code modifications.

% === 4.3 会话状态管理 ===
\subsection{Session State Management}

% 中文注释：Session的KV存储实现
Session implements a KV store pattern with topic-based result sharing:

\begin{itemize}[leftmargin=1.5em]
    \item \textbf{Plan storage:} Current Plan accessible to all agents.
    \item \textbf{Topic registry:} Results keyed by topic name (from PlanStep.produces).
    \item \textbf{Metadata tracking:} Execution timestamps, agent identifiers, error logs.
\end{itemize}

Topics enable data flow between dependent steps: a step producing topic ``schema\_info'' stores results that downstream steps consuming ``schema\_info'' can retrieve.

% === 4.4 安全机制 ===
\subsection{Security Mechanisms}

% 中文注释：安全防护措施
The system implements multiple security layers:

\paragraph{Path Validation.}
File tools validate paths against the designated context directory, preventing directory traversal attacks. Paths are normalized and compared to the canonical context path.

\paragraph{SQL Injection Prevention.}
SQL execution tools parse queries to verify SELECT-only operations. Keywords like INSERT, UPDATE, DELETE, DROP are rejected before execution.

\paragraph{Execution Timeout.}
Python execution imposes a 30-second timeout via signal-based interruption, preventing runaway computations from blocking the system.

\paragraph{Tool Registry.}
Only registered tools are invocable; arbitrary function execution is prohibited.


% === 第五章：实验分析 ===
\section{Experimental Analysis}

% 中文注释：实验分析章节概述
We evaluate ElianDataAgent on representative data analysis scenarios to assess planning quality, execution efficiency, and replanning effectiveness.

% === 5.1 实验设置 ===
\subsection{Experimental Setup}

% 中文注释：实验环境和配置
\paragraph{Environment.}
Experiments run on a standard configuration:
\begin{itemize}[leftmargin=1.5em]
    \item LLM Backend: Qwen2.5-72B \cite{qwen2.5} or GPT-4o \cite{hurst2024gpt}.
    \item Data Sources: Mix of CSV files, JSON records, SQLite databases.
    \item Execution: DAG-enabled mode with Hybrid strategy.
\end{itemize}

\paragraph{Tasks.}
% 中文注释：测试任务类型
We define task categories representing common data analysis patterns:
\begin{itemize}[leftmargin=1.5em]
    \item \textbf{Single-source exploration:} Query a single CSV file for statistics or patterns.
    \item \textbf{Multi-source integration:} Combine data across multiple tables or files.
    \item \textbf{SQL query formulation:} Construct and execute complex SQL queries.
    \item \textbf{Statistical computation:} Perform aggregations, distributions, or correlations.
\end{itemize}

% === 5.2 规划质量 ===
\subsection{Planning Quality}

% 中文注释：规划质量评估
We assess Planner output quality through manual inspection of generated Plans:

\begin{table}[t]
\centering
\caption{Planning quality metrics across task types.}
\label{tab:plan-quality}
\begin{tabular}{lccc}
\toprule
Task Type & Avg Steps & Dependency Accuracy & Mode Selection \\
\midrule
Single-source & 2.3 & 95\% & Sequential (80\%) \\
Multi-source & 4.7 & 92\% & Hybrid (85\%) \\
SQL formulation & 3.1 & 88\% & Sequential (60\%) \\
Statistical & 3.8 & 90\% & Hybrid (75\%) \\
\bottomrule
\end{tabular}
\end{table}

Dependency accuracy measures whether declared dependencies correctly reflect logical operation ordering. Mode selection indicates appropriate strategy choices (Sequential for simple tasks, Hybrid for parallelizable tasks).

% === 5.3 执行效率 ===
\subsection{Execution Efficiency}

% 中文注释：执行效率对比
We compare execution times between Sequential and DAG (Hybrid) modes:

\begin{table}[t]
\centering
\caption{Execution efficiency comparison (seconds).}
\label{tab:exec-efficiency}
\begin{tabular}{lccc}
\toprule
Task Type & Sequential & DAG Hybrid & Improvement \\
\midrule
Multi-source integration & 45.2 & 28.7 & 37\% \\
Statistical computation & 32.1 & 21.3 & 34\% \\
Complex SQL query & 18.5 & 12.2 & 35\% \\
\bottomrule
\end{tabular}
\end{table}

DAG-enabled execution shows consistent efficiency improvements (30-40\%) for tasks with parallelizable operations, validating the value of dependency-aware concurrent execution.

% === 5.4 重规划有效性 ===
\subsection{Replanning Effectiveness}

% 中文注释：重规划场景分析
We examine replanning behavior in scenarios requiring plan modification:

\begin{itemize}[leftmargin=1.5em]
    \item \textbf{Missing column detection:} Executor discovers a referenced column does not exist; Replanner adds schema inspection step and modifies query formulation.
    \item \textbf{Unexpected data format:} CSV parsing reveals inconsistent encoding; Replanner adds preprocessing step for encoding normalization.
    \item \textbf{Query timeout:} Complex SQL exceeds execution limits; Replanner suggests query decomposition into simpler subqueries.
\end{itemize}

Replanning success rate (tasks completed after replanning) reaches 82\%, demonstrating effective adaptation to intermediate discoveries.

% === 5.5 案例研究 ===
\subsection{Case Study}

% 中文注释：案例分析展示
We present a representative case illustrating the complete workflow:

\textbf{Task:} ``Analyze the relationship between customer demographics and purchase amounts across three CSV files.''

\begin{enumerate}[leftmargin=1.5em]
    \item \textbf{Planning:} Planner generates 5-step Plan:
    \begin{itemize}[leftmargin=1em]
        \item $s_1$: Explore file structures (produces: ``schema\_info'')
        \item $s_2$: Load demographics data (depends: $s_1$; produces: ``demo\_data'')
        \item $s_3$: Load purchase data (depends: $s_1$; produces: ``purchase\_data'')
        \item $s_4$: Join datasets and compute correlations (depends: $s_2, s_3$; produces: ``analysis\_result'')
        \item $s_5$: Synthesize answer (depends: $s_4$; produces: ``final\_answer'')
    \end{itemize}

    \item \textbf{Execution:} DAGScheduler selects Hybrid strategy.
    \begin{itemize}[leftmargin=1em]
        \item Layer 0: $s_1$ executes alone.
        \item Layer 1: $s_2, s_3$ execute concurrently (no inter-dependency).
        \item Layer 2: $s_4$ executes after $s_2, s_3$ complete.
        \item Layer 3: $s_5$ synthesizes final answer.
    \end{itemize}

    \item \textbf{Replanning:} During $s_3$ execution, missing customer IDs detected; Replanner adds preprocessing step for ID normalization; subsequent steps proceed with adjusted Plan.
\end{enumerate}

This case demonstrates the framework's ability to decompose complex tasks, parallelize independent operations, and adapt to execution discoveries.


% === 第六章：结论 ===
\section{Conclusions and Future Work}

% 中文注释：结论总结
ElianDataAgent provides a practical multi-agent framework for data analysis that integrates Plan-Execute-Replan orchestration with DAG-enhanced parallel execution. The three-agent architecture (Planner, Executor, Replanner) enables strategic task decomposition, step-wise tool reasoning, and dynamic plan adaptation, while the DAG scheduler maximizes execution efficiency through dependency-aware concurrency.

% 中文注释：技术贡献总结
Key technical contributions include:
\begin{itemize}[leftmargin=*, align=left]
    \item A comprehensive type system for inter-agent communication with Session-based state management.
    \item Three execution strategies (Sequential, Parallel, Hybrid) with Kahn's algorithm-based topological ordering.
    \item A secure tool registry with validated data manipulation capabilities.
    \item Replanning mechanisms for adaptive plan modification.
    \item OpenAI API compatibility supporting both 0.x and 1.x versions.
\end{itemize}

% 中文注释：未来工作展望
Future work directions include:
\begin{itemize}[leftmargin=*, align=left]
    \item \textbf{Agent memory enhancement:} Integrating long-term memory mechanisms \cite{zhong2024memorybank,packer2023memgpt} for cross-session knowledge retention.
    \item \textbf{Tool library expansion:} Extending the tool registry for specialized domains (financial analysis, scientific computing).
    \item \textbf{Benchmark integration:} Systematic evaluation on established agent benchmarks \cite{liu2023agentbench,zheng2025lifelongagentbench}.
    \item \textbf{Multi-model support:} Enabling heterogeneous agent configurations where different agents use specialized LLMs.
\end{itemize}

ElianDataAgent establishes a foundation for building intelligent data analysis systems that combine strategic planning with efficient execution, pointing toward scalable and adaptive agent architectures for complex analytical tasks.


% === 参考文献 ===
\bibliographystyle{unsrt}
\bibliography{references}

\end{document}